\section{国内外研究现状}
\subsection{国内研究现状}
国内幻灯片制作工具的研究主要集中在两个方向：传统PPT工具的功能拓展与基于Web的轻量化解决方案。早期研究聚焦于Microsoft PowerPoint等传统工具的本土化模板开发与插件适配，解决了通用场景下的基础制作需求，但操作复杂度较高。

随着Web技术与Markdown语言的普及，国内开始研究基于HTML5的幻灯片制作工具。基于HTML5的幻灯片设计及实现智能手机播放的研究\cite{HLJK201802015}探讨了如何利用现代Web技术制作跨平台幻灯片。HTML5互动微课开发与应用\cite{ZJGZ201603011}和基于HTML5技术的在线课程的设计与制作\cite{XXDL201613117}的研究进一步证明了Web技术在教育演示领域的应用潜力。这些研究为基于Web的幻灯片工具提供了重要基础，但尚未深入探讨团队协作功能。

在权限管理方面，国内已有较为成熟的研究成果。基于角色的权限管理系统的设计与实现\cite{zhang2013rbac}为权限控制提供了系统设计思路；基于角色访问控制（RBAC）的Web应用研究\cite{wu2002rbacweb}探讨了RBAC模型在Web环境中的实践；基于RBAC模型云资源管理系统访问控制权限的设计与实施\cite{wei2025rbaccloud}则针对云环境提供了解决方案。这些研究为本平台权限系统的设计提供了理论基础。

在文档协作领域，基于目录节点的WEB文档协作编辑系统分析与设计\cite{tan2012web}和基于用户体验的CSCW在线文档平台交互设计研究\cite{li2020cscw}为多人协作编辑提供了设计参考。基于FileNet的文档共享系统的设计与实现\cite{li2015filenet}和基于Java Web组件技术的软件项目文档管理系统\cite{zhang2024javaweb}则提供了云端资源管理的技术思路。

在插件化架构方面，在线文档中基于微前端的容器化插件系统的设计与实现\cite{cai2022microfrontend}为可扩展的插件系统设计提供了重要参考。这些研究为本平台的模块化设计奠定了坚实基础。

\subsection{国外研究现状}
国外在轻量化幻灯片工具领域的研究起步较早，形成了以Reveal.js、Marp、Slidev等为代表的成熟生态。其中，Slidev (slide + dev, /slaɪdɪv/) 是一个基于Web的现代幻灯片制作工具，它以Markdown的形式专注于编写幻灯片内容，能够制作出具有交互式演示功能的、高度可自定义的幻灯片。基于HTML5的跨平台移动应用关键技术的研究与实现\cite{GYKJ201303028}为这类基于Web的工具提供了重要的技术支撑。

在实时协作技术方面，国外研究已形成了较为完善的理论体系。WebSocket在Web实时通信领域的研究\cite{DNZS201028018}为实时数据传输提供了解决方案。研究集中在增量同步、操作转换等算法上，以解决多端编辑的冲突问题，确保数据一致性。这些技术为多人协作编辑功能提供了理论依据。

在权限控制方面，国外学者对角色访问控制进行了深入研究。Employing UML and OCL for designing and analysing role-based access control\cite{kuhlmann2013uml}探讨了使用形式化方法设计RBAC系统；Feasibility study of software reengineering towards role-based access control\cite{li2011reengineering}则研究了向RBAC系统迁移的可行性。这些研究为权限系统的设计提供了理论指导。

在版本控制方面，Git版本控制系统在软件开发中的应用与实践\cite{SXDS202505026}和版本控制工具在软件开发项目管理中的应用\cite{XMGJ202006028}的研究展示了版本控制在团队协作中的重要性，但同时也指出了传统版本控制工具对非技术用户的复杂性挑战。

\section{研究成果}
现有研究成果可分为三类：传统可视化工具、基于Markdown的轻量化工具以及协作相关技术。

传统可视化工具以Microsoft PowerPoint和Apple Keynote为代表，采用所见即所得编辑模式，提供丰富的模板和可视化操作界面。这类工具的优点是直观易用，适合非技术用户快速创建基础演示文稿。然而，它们存在明显局限性：操作相对复杂，不支持Markdown等轻量级标记语言，技术内容渲染依赖第三方插件，协作功能仅限于文件共享，缺乏细粒度的版本控制能力。

基于Markdown的轻量化工具方面，Reveal.js是一个开源HTML幻灯片框架，支持使用Markdown创建可定制化幻灯片，并允许通过HTML和CSS实现高度自定义。其优点是扩展性强、功能丰富。Marp专注于简洁性和可移植性，提供极简的Markdown转幻灯片流程，支持多种输出格式。其优点是轻量、易上手，但功能相对有限。

Slidev整合了上述工具的优势，是一个专门为开发者设计的现代幻灯片制作工具。Slidev的核心创新在于其基于扩展Markdown的插件化架构，在保持Markdown简洁性的同时，通过丰富的插件系统支持各种高级功能。Slidev具有四个显著特点：开发者友好（原生支持代码片段、终端模拟等开发元素）、基于现代Web技术栈（充分利用Vue、Vite等框架）、高度可定制（允许使用Web技术自定义样式与交互）、开源且有活跃的社区支持。相较于Reveal.js，Slidev更加简洁且支持Vue组件；相较于Marp，它在保持易用性的同时提供了更丰富的功能。然而，Slidev当前缺乏对团队协作的原生支持。

在多人协作技术方面，实时协同编辑算法主要包括增量同步和操作转换两大类。增量同步算法通过传输编辑差异实现高效数据传输，实时性强但对冲突检测精度要求高；操作转换算法通过数学方法解决多端操作时序冲突。

在权限管理领域，RBAC模型已在各类系统中得到广泛应用，形成了较为完善的理论体系。然而，这些研究尚未与Slidev的协作场景深度结合，特别是在权限粒度设计方面有待探索。

在插件化架构方面，基于微前端的容器化插件系统研究成果，为实现功能的模块化扩展提供了技术方案，但需要针对Slidev的特殊语法进行适应性改造。

\section{进展情况和发展趋势}
当前，轻量化幻灯片工具领域的发展呈现以下特征：

技术生态方面，Slidev已经建立了较为完善的插件化体系，为集成协作功能和AI辅助能力奠定了良好基础。协作技术方面，实时协同编辑算法在通用文档工具中已相对成熟，为构建Slidev的多人协作功能提供了可直接借鉴的方案。权限管理方面，基于RBAC的权限模型已在各类系统中得到充分验证，可应用于Slidev协作平台的权限系统设计。AI集成方面，大语言模型的发展为幻灯片内容的智能生成提供了技术支持，提示词工程的进步使得针对Slidev语法的专用模板成为可能。

未来发展趋势呈现三个方向：

第一是协作化趋势。随着团队协作需求的增长，轻量化幻灯片工具将从单机版向多人实时协同平台演进。基于RBAC的权限管理机制和优化的协作交互逻辑将成为平台基础能力，以降低非技术团队的协作门槛。

第二是智能化趋势。AI技术将从单纯内容生成向语法适配、格式优化、模板推荐等全流程辅助延伸。通过结合插件化架构，AI功能可以模块化扩展，提升生成内容与Slidev语法的适配性。

第三是轻量化趋势。在保证功能完整性的前提下，专业协作工具的操作流程将进一步简化，以适应特定场景的团队使用习惯。

本研究紧密贴合上述发展趋势，通过设计轻量级冲突检测机制、简易版本控制系统和智能语法适配模块，解决Slidev协作能力缺失的核心问题，具有理论价值和实践意义。

\section{存在的问题}
通过对现有研究成果的梳理分析，本研究发现当前Slidev生态在团队协作方面存在以下五个突出问题：

\begin{enumerate}[label=(\arabic*)]
    \item Slidev框架缺乏团队协作核心能力。现有Slidev工具仅支持单机编辑模式，多人协作需要依赖外部文件传输工具，导致版本混乱、权限管理失控。虽然RBAC权限管理方案为权限体系设计提供了参考，但这些方案尚未针对Slidev的协作场景进行专门适配，难以满足团队化制作需求。
    
    \item 云端资源管理与Slidev的适配性不足。现有云端存储方案大多采用通用路径格式，不符合Slidev特有的资源引用规范。虽然文档共享系统和文档管理系统在资源管理方面有成熟设计，但这些系统未针对Slidev的资源引用特性进行适配，导致云端存储的图片、模板等资源无法被Slidev项目直接调用。
    
    \item AI生成内容与Slidev语法的适配性差。通用大语言模型生成的文本内容通常无法直接转换为符合Slidev规范的代码块、组件指令、动画脚本（如v-click）等特殊语法元素。尽管微前端插件化架构为实现AI功能模块扩展提供了技术可能，但现有研究尚未针对Slidev的特殊语法设计专门的生成与转换逻辑。
    
    \item 版本控制工具的复杂性影响非技术团队的协作效率。专业的版本控制工具（如Git）操作相对复杂，包含分支管理、合并冲突等高级概念，而Slidev项目基于Markdown的文件结构相对简单，通常不需要复杂的版本操作。对于非技术背景的团队成员而言，现有的版本控制系统学习成本较高，影响了协作体验。
    
    \item 云资源权限管控与Slidev协作场景脱节。虽然基于RBAC的云资源管理权限设计在理论和技术上较为成熟，但现有方案未充分考虑Slidev轻量化协作的特点，权限粒度的设置与团队协作的实际需求存在偏差。
\end{enumerate}

\section{主要参考文献}
% 模板会自动根据正文中的\cite引用生成参考文献列表
% 已完整引用以下16篇文献，每篇仅引用一次：
% HLJK201802015, ZJGZ201603011, XXDL201613117, zhang2013rbac, wu2002rbacweb, 
% wei2025rbaccloud, tan2012web, li2020cscw, li2015filenet, zhang2024javaweb, 
% cai2022microfrontend, GYKJ201303028, kuhlmann2013uml, li2011reengineering, 
% DNZS201028018, SXDS202505026, XMGJ202006028